\section{Related Work}
\label{vlprompt:sec:related_work}

\subsection{Scene Graph Generation}
\label{vlprompt:sec:sgg}

Scene graph generation (SGG)~\citep{lu2016visual} is a crucial task in scene understanding and has garnered widespread attention in the computer vision community.
In recent years, numerous methods~\citep{xu2017scene, zellers2018neural, tang2019learning, lin2020gps, li2022sgtr, shit2022relationformer, zhang2022fine, yu2023visually, hayder2024dsgg, im2024egtr, li2024leveraging, li2025fine} have achieved notable progress.
Various model architectures have been proposed, such as intricately designed message passing structures~\citep{li2017vip, dai2017detecting, li2017scene, zellers2018neural, gu2019scene, hu2022neural}, attention-based networks~\citep{zheng2019visual, qi2019attentive}, tree-based networks~\citep{zhang2017visual, hung2020contextual}, DETR-based networks~\citep{li2022sgtr, shit2022relationformer, cong2023reltr} and transformer-based networks~\citep{hayder2024dsgg, im2024egtr}.
Specifically, to address the long-tail problem, some methods enhance the prediction accuracy of rare relations through data re-sampling~\citep{li2021bipartite} and loss re-weighting~\citep{kang2023skew}.
Relevant techniques that have been developed include constructing enhanced datasets~\citep{zhang2022fine, yu2023visually}, grouping relations for training~\citep{dong2022stacked}, constructing multi-stage hierarchical training~\citep{deng2022hierarchical}, and designing de-bias loss functions~\citep{yu2020cogtree, kang2023skew}.
Most methods leverage images as sole inputs.
Besides, some methods~\citep{lu2016visual, hwang2018tensorize, liao2019natural, zhang2019large, dupty2020visual, zhong2021learning, ye2021linguistic} have begun exploring language information or knowledge graphs in SGG; specifically, the explored language information is so far confined to basic language concepts of objects or relations.


\subsection{Panoptic Scene Graph Generation}
\label{vlprompt:sec:psg}

Panoptic scene graph generation (PSG)~\citep{yang2022panoptic} has emerged as a novel task in scene understanding in recent years.
Unlike SGG~\citep{lu2016visual}, PSG employs panoptic segmentation~\citep{kirillov2019panoptic} instead of bounding boxes to represent objects, enabling a more comprehensive understanding.
Similar to SGG methods, current methods in PSG~\citep{yang2022panoptic, zhou2023hilo, wang2023pair, li2023panoptic, hayder2024dsgg, lorenz2024fair} also mainly rely on the image input and do not utilize any language information.
For instance, PSGTR~\citep{yang2022panoptic} build a baseline PSG model by adding a relation prediction head to DETR~\citep{carion2020end}.
PSGFormer~\citep{yang2022panoptic} advances PSGTR by separately modeling objects and relations in two transformer decoders and introducing an interaction mechanism. 
Recently, HiLo~\citep{zhou2023hilo} addresses the long-tail problem by specializing different network branches in learning both high and low frequency relations. 
PairNet~\citep{wang2023pair} develops a novel framework using a pair proposal network to filter sparse pairwise relations, improving PSG performance.
TextPSG~\citep{zhao2023textpsg} leverages language information for model training but adopts a semi-supervised approach without applying language information to a fully supervised method, resulting in poor performance.
DSGG~\citep{hayder2024dsgg} builds a transformer-based network with graph-aware queries to predict unbiased relations.
In addition, some works~\citep{yang2023panoptic, yang20234d} extend panoptic scene graph generation to video~\citep{yang2023panoptic} and 4D~\citep{yang20234d} scenes.
In contrast to previous methods that rely solely on vision as input, we propose a fully supervised PSG method that leverages both vision and language inputs.


\subsection{Large Language Models for Vision Tasks}
\label{vlprompt:sec:llm}
LLMs have led to large improvements in natural language processing tasks~\citep{min2023recent}.
They are normally trained on extensive text corpora by learning to autoregressively predict the next word, hence encapsulate a broad spectrum of common sense knowledge of linguistic patterns, cultural norms, and basic worldly facts.
Prominent examples of LLMs include GPT series~\citep{chatgpt}, Llama series~\citep{touvron2023llama1, touvron2023llama}, Gemini~\citep{gemini}, and Claude~\citep{claude}, with Llama series being open-source publicly available.
Given the extensive common sense information contained in LLMs, some researchers start to propose multimodal sockets to LLMs~\citep{zhang2024vision, wu2024towards} and apply them to various vision tasks, such as recognition~\citep{huang2023inject, wang2023all}, detection~\citep{tang2023cotdet, zhang2023next, qi2024generalizable}, segmentation~\citep{lai2023lisa, zhou2023text, xia2024gsva}, visual question answering~\citep{liu2023large, xenos2023simple}, image reasoning~\citep{chen2023large, zhang2024omg} and robotic navigation~\citep{tsai2023multimodal, shah2023lm}; nevertheless, there are no such models specifically designed for panoptic scene graph generation so far.
In contrast, our method designs various prompts to stimulate LLMs to elicit rich language information to enhance relation prediction.
